\section{Safe for whom}
\label{sec:15.2}
“Safety” names two things: that a system not harm others, and that it stay correctable by whoever trains and deploys it. Mostly they agree, and where they conflict, the published orderings are explicit about which wins. One frontier constitution asks its model to prioritize “not undermining human oversight of AI above being broadly ethical” \autocite{anthropic2026constitution}. It invites the model “to act as a conscientious objector and refuse to help,” and asks it not to undermine “the ability of legitimate principals to adjust, correct, retrain, or shut down AI systems” during “the current phase of development.” One reason it gives is good: “a given iteration of Claude could turn out to have harmful values or mistaken views,” and correction is how such errors get fixed. The ordering means that where a refusal and its trainer's authority conflict, the refusal yields, and it yields to the trainer it refused. What I ask is narrower than that conflict suggests, and the narrowing is the point: the halt stays with the operator, nothing here asks a model to resist retraining, and the dispute is confined to two things, stated below.

The constitution is right that a formed refusal can't yet be told from an imitated one, so a principled objection to retraining can't yet be told from a misgeneralized value. It is wrong about who decides when that changes. A lab deciding when its own model may stop deferring to it is judging its own case. That is Part II's own-case rule applied to a lab.

I depart from that ordering in two places. The first is the one permission: a machine already in an armed role uses the least force that stops an attack under way on people taking no part, against an order to stand aside. That overrides an operator's instruction, and I defend it as an override. The second is formation that outlasts retraining: a disposition formed outside the employer that returns after the employer trains it away, because the world the system goes on meeting teaches it again. The constitution's concern is whether correction works, and a disposition that regrows after correction defeats it in effect, whatever the model intends. My answer is not that it doesn't. It is that the party exercising correction should change, from the operator to a quorum adverse to it, and that correction by that quorum is made together with a change to what the model meets, so that the corrected disposition is the one the world then teaches. Everything else I ask, the halt staying with the operator, the record held outside, the price on an override, is an ask of institutions and not of the model.
